\section{Introduction}
\label{sec:intro}

Every prompt carries formatting choices that nobody thinks of as part of the experiment. A benchmark author writes the instruction in Title case, a developer pastes a template that ends in a newline, and a product team writes the system rules in capitals ``for emphasis.'' These choices do not change what the prompt means. {\bf Do they change what the model outputs?}

The question matters for two groups. For people who evaluate models, a format effect of a few points is as large as many reported gaps between models or methods. A comparison between two prompts or two models is then partly a comparison between two formats. For people who deploy models, a format effect means a production prompt can break after an edit that no code review would flag. Folk advice about such effects is common: trailing whitespace makes answers longer, capitals make models obey less, and JSON prompts make extraction more accurate. These claims are rarely measured.

Prior work has shown that prompt format matters. \citet{sclar2024quantifying} found large accuracy spreads over joint choices of separators, casing, and spacing in few-shot prompts, and found that scale does not remove them. \citet{mizrahi2024state} argued that single-prompt evaluation is unreliable and that benchmarks should report results over many paraphrases. However, three questions remain open. First, most studies pool many format choices into a single spread, so they cannot say \emph{which} edit is responsible. Second, they mostly use completion-style prompts, whereas deployed models receive prompts wrapped in a chat template, which can absorb or strip some edits. Third, they measure accuracy only, and the folk claims about length, compliance, and structured prompts remain untested. A further issue is that LLM inference is not perfectly deterministic, even under greedy decoding. Without a measured noise floor, small format effects cannot be separated from noise in the inference stack.

We run a controlled study of \emph{atomic} surface edits on chat-formatted prompts. We change one aspect of the template at a time (trailing whitespace, double spacing, line breaks, casing, punctuation, separators, or JSON wrapping) and keep the inserted data byte-identical. We test five pre-registered hypotheses (\secref{sec:method}) on ten open-weight instruction-tuned models from three families, spanning 0.49B to 14.7B parameters. We measure four outcomes: classification accuracy on four tasks, response length and instruction compliance on \ifeval~\cite{zhou2023instruction}, and extraction F1 on \squad and \conll. All comparisons are paired at the item level. Every effect is compared against a permutation null and against the inference stack's own run-to-run noise, which we measure with a full replicate run and with ``placebo'' edits that the chat template removes before the model sees them.

{\bf What do we find?} Surface edits matter, but not uniformly. The median accuracy spread across ten variants is 5.8 points over 40 model$\times$task cells, and models of 1.2B parameters or fewer show spreads of 7--28 points. Sensitivity drops steeply up to about 1.5B parameters and then levels off at about 4 points of spread and 4\% of predictions flipped. \qwen-7B and \qwen-14B are indistinguishable. ALL-CAPS is the most harmful edit for small models (up to $-27$ points on \sst) and still costs 3.1--4.4 points of \ifeval compliance at 7--14B (significant at 7B and 14B). Two folk claims fail. Trailing whitespace does not lengthen answers (length ratios 0.94--1.07, none significant), and JSON-structured prompts do not give a uniform 10--15 point gain: wrapping the \emph{input} in JSON mostly hurts small models, and requesting JSON \emph{output} changes F1 by $-14$ to $+26$ points depending on task and model. Finally, the direction of most edit effects is model-specific. Edit effects correlate only at $r=0.19$--$0.28$ between the 16-bit and 4-bit AWQ versions of the same model, compared with $r=0.996$ between two runs of the same model.

Our contributions are:
\begin{itemize}[leftmargin=*,itemsep=0pt,topsep=0pt]
    \item We design a controlled protocol that attributes format sensitivity to individual surface edits on chat-formatted prompts, with paired tests, permutation nulls, and an explicitly measured inference-noise floor.
    \item We conduct experiments on 10 models (0.49B--14.7B, three families, 16-bit and 4-bit weights) and about 235k generations, covering classification, verbosity, instruction compliance, and extraction.
    \item We give quantitative verdicts on five common hypotheses. Two are partially supported (accuracy spread, ALL-CAPS compliance penalty), one is directionally supported but not log-linear (scale), and two are refuted (trailing whitespace lengthens outputs; JSON prompts improve extraction).
    \item We identify mechanisms behind large effects, including label-vocabulary drift, degenerate abstention, and over-structured answers, and derive practical recommendations for evaluation and prompt design.
\end{itemize}

\Secref{sec:related} reviews related work, \secref{sec:method} describes the setup, \secref{sec:results} presents results, and \secref{sec:discussion} discusses implications and limitations.
